\documentclass[12pt,a4paper]{article}

% --- PAQUETES DE IDIOMA Y FORMATO ---
\usepackage[utf8]{inputenc}
\usepackage[spanish,es-tabla]{babel}
\usepackage{geometry}
\geometry{top=2.5cm, bottom=2.5cm, left=3cm, right=2.5cm}
\usepackage{graphicx}
\usepackage{booktabs}
\usepackage{tabularx}
\usepackage{longtable}
\usepackage{array}
\usepackage{colortbl}
\usepackage{xcolor}
\usepackage{hyperref}
\usepackage{enumitem}
\usepackage{fancyhdr}
\usepackage{microtype}
\usepackage{tcolorbox}
\usepackage{cite}

% --- DEFINICIÓN DE COLORES INSTITUCIONALES Y TÉCNICOS ---
\definecolor{unsaacBlue}{RGB}{0, 51, 102}
\definecolor{unsaacGold}{RGB}{179, 134, 0}
\definecolor{tableHeader}{RGB}{235, 242, 250}
\definecolor{tableStripe}{RGB}{248, 250, 253}
\definecolor{tierAStar}{RGB}{220, 53, 69}
\definecolor{tierA}{RGB}{40, 167, 69}
\definecolor{tierB}{RGB}{23, 162, 184}
\definecolor{darkgray}{RGB}{50, 50, 50}

% --- CONFIGURACIÓN DE HIPERVÍNCULOS ---
\hypersetup{
    colorlinks=true,
    linkcolor=unsaacBlue,
    citecolor=unsaacBlue,
    urlcolor=unsaacBlue,
    pdfauthor={Grupo 4 - UNSAAC},
    pdftitle={Investigacion y Clasificacion de Conferencias en Desarrollo de Software}
}

% --- ENCABEZADOS Y PIES DE PÁGINA ---
\pagestyle{fancy}
\fancyhf{}
\fancyhead[L]{\small\color{gray}UNSAAC -- Ingeniería de Software I}
\fancyhead[R]{\small\color{gray}Principales Conferencias de Software}
\fancyfoot[C]{\thepage}
\renewcommand{\headrulewidth}{0.4pt}

% ==============================================================================
% COMIENZO DEL DOCUMENTO
% ==============================================================================
\begin{document}

% --- PORTADA ---
\begin{titlepage}
    \centering
    \vfill
    {\large\bfseries UNIVERSIDAD NACIONAL DE SAN ANTONIO ABAD DEL CUSCO} \\[0.4cm]
    {\large FACULTAD DE INGENIERÍA ELÉCTRICA, ELECTRÓNICA, INFORMÁTICA Y MECÁNICA} \\[0.4cm]
    {\large ESCUELA PROFESIONAL DE INGENIERÍA INFORMÁTICA Y DE SISTEMAS} \\[1.5cm]
    
    \includegraphics[width=0.35\textwidth]{unsaac.jpg} \\[1.5cm] % Nota: Cambiar por logo local si es necesario, o Overleaf descargará automáticamente.
    
    {\Large\bfseries ASIGNATURA: INGENIERÍA DE SOFTWARE I} \\[0.6cm]
    {\large\bfseries TRABAJO: INVESTIGACIÓN Y CLASIFICACIÓN DE LAS PRINCIPALES CONFERENCIAS SOBRE DESARROLLO DE SOFTWARE} \\[0.6cm]
    {\large\bfseries RELEVANCIA, RANKINGS (CORE / GOOGLE SCHOLAR) Y MATRIZ DE CALIFICACIÓN} \\[1.5cm]
    
    \textbf{DOCENTE RESPONSABLE:} \\
    Mgt. Luz Indira Ticona Felix \\[1.0cm]
    
    \textbf{INTEGRANTES / DESARROLLADORES (GRUPO 4):} \\
    ALVAREZ CATUNTA, Angel Ismael - 221945 \\
    CHILLIHUANI HUAMAN, Rivaldo Franco - 221445 \\
    SUPA CUSIPAUCAR, Yeferson - 220553 \\
    VILLAVICENCIO SEGUIL, Edu Piero - 220555 \\[0.5cm]
    
    {\large Cusco - Perú} \\
    {\large 2026}
    \vfill
\end{titlepage}

\newpage
\tableofcontents
\newpage

% ==============================================================================
% 1. INTRODUCCIÓN
% ==============================================================================
\section{Introducción}

\subsection{El Papel Crítico de las Conferencias en la Ciencia de la Computación}
En la mayoría de disciplinas académicas tradicionales (como la medicina, la física teórica o la química), las revistas científicas indexadas (\textit{journals}) representan la vía por excelencia para la comunicación de descubrimientos de alto impacto. Sin embargo, en las Ciencias de la Computación y, en particular, en la \textbf{Ingeniería y Desarrollo de Software}, existe un fenómeno metodológico e histórico particular: \textbf{las conferencias internacionales arbitradas constituyen el canal primario, más riguroso y prestigioso para la difusión de avances científicos y tecnológicos}.

Esta dinámica se debe a tres factores primordiales:
\begin{enumerate}[leftmargin=*]
    \item \textbf{Celeridad del Ciclo de Innovación:} La industria del software y sus paradigmas evolucionan a un ritmo vertiginoso. Los tiempos habituales de revisión editorial de una revista científica (que pueden superar los doce o dieciocho meses) resultan incompatibles con la rápida obsolescencia de las herramientas, lenguajes y arquitecturas de software.
    \item \textbf{Rigor del Proceso de Arbitraje:} Las conferencias de primer nivel implementan revisiones por pares a doble ciego (\textit{double-blind review}), etapas de refutación formal (\textit{author rebuttal phase}) y comités de evaluación de artefactos de software (\textit{Artifact Evaluation Committees}), garantizando la reproducibilidad del código, datos y experimentos.
    \item \textbf{Tasas de Aceptación Extremadamente Selectivas:} Las conferencias élite aceptan con frecuencia menos del 20\% de los artículos postulados, superando en selectividad a gran parte de las revistas indexadas de nivel medio.
\end{enumerate}

\subsection{Objetivos del Informe}
El presente trabajo tiene por propósito:
\begin{itemize}[leftmargin=*]
    \item Investigar y caracterizar las conferencias más influyentes en el ámbito del desarrollo y la ingeniería de software a escala global.
    \item Ordenar y clasificar dichas conferencias de acuerdo con los sistemas de calificación (\textit{qualification}) reconocidos por la comunidad científica internacional (especialmente el ranking \textbf{CORE}, métricas bibliométricas de \textbf{Google Scholar Metrics} y \textbf{CSRankings}).
    \item Ofrecer un análisis comparativo entre los espacios de investigación científica y las conferencias de la industria profesional, subrayando las tendencias emergentes que impactan la formación de los futuros ingenieros informáticos.
\end{itemize}

% ==============================================================================
% 2. MARCOS DE CALIFICACIÓN Y EVALUACIÓN DE RELEVANCIA
% ==============================================================================
\section{Sistemas de Evaluación y Calificación de Conferencias}

Para determinar la relevancia e impacto de una conferencia científica en desarrollo de software se recurre a marcos formales estandarizados:

\subsection{El Ranking CORE (Computing Research and Education Association of Australasia)}
El ranking CORE es considerado el estándar de oro para la categorización cualitativa de conferencias en computación en todo el mundo. Las clasificaciones se determinan tras auditorías de comités internacionales de expertos que evalúan el prestigio de los comités de programa, la tasa de aceptación, la longevidad del evento y el impacto de sus publicaciones:
\begin{itemize}[leftmargin=*]
    \item \textbf{Nivel A* (Flagship / Élite Global):} Representa típicamente el 4\% al 7\% superior de todas las conferencias del área. Son los foros de máxima reputación académica, donde se definen las tendencias teóricas y metodológicas de la disciplina.
    \item \textbf{Nivel A (Excelente / Alta Calidad):} Agrupa a conferencias de gran renombre internacional (aproximadamente el 15\% superior), con una selectividad rigurosa y un impacto consolidado en subdisciplinas especializadas.
    \item \textbf{Nivel B (Buena Calidad / Reconocidas):} Conferencias de buen nivel técnico con presencia internacional o regional sólida, con tasas de aceptación moderadas.
    \item \textbf{Nivel C (Iniciación / Ámbito Regional):} Foros arbitrados con estándares de selectividad menores, orientados a comunidades locales o difusión preliminar de trabajos.
\end{itemize}

\subsection{Métricas Bibliométricas de Google Scholar ($h5\text{-index}$ y $h5\text{-median}$)}
Google Scholar Metrics evalúa el impacto cuantitativo de las publicaciones en los últimos cinco años completos:
\begin{itemize}[leftmargin=*]
    \item \textbf{Índice $h5$:} Número $h$ de artículos publicados en los últimos 5 años que han recibido al menos $h$ citas cada uno. Refleja la consistencia del impacto del evento.
    \item \textbf{Mediana $h5$:} La mediana del número de citas recibidas por los $h$ artículos que conforman el índice $h5$. Refleja la profundidad del impacto bibliográfico.
\end{itemize}

\subsection{CSRankings y Rating GII-GRIN-SCIE}
\begin{itemize}[leftmargin=*]
    \item \textbf{CSRankings:} Métricas desarrolladas por investigadores académicos internacionales para medir la productividad institucional basada exclusivamente en publicaciones aceptadas en los foros más selectivos (en Ingeniería de Software: ICSE, FSE y ASE).
    \item \textbf{GII-GRIN-SCIE (GGS):} Sistema promovido por las sociedades científicas de computación de Italia y España que clasifica las conferencias en Clase 1, Clase 2 y Clase 3 según su visibilidad y factor de impacto normalizado.
\end{itemize}

% ==============================================================================
% 3. CLASIFICACIÓN JERÁRQUICA DE LAS CONFERENCIAS
% ==============================================================================
\section{Principales Conferencias Ordenadas por Relevancia}

A continuación, se presentan las conferencias científicas y técnicas más influyentes en desarrollo de software, estructuradas de forma descendente según su calificación oficial:

% --- 3.1 NIVEL CORE A* ---
\subsection{Nivel 1: Categoría CORE A* (Flagship / Élite Mundial)}

\begin{tcolorbox}[colback=white,colframe=tierAStar,title={\textbf{1. ICSE -- International Conference on Software Engineering}},fonttitle=\bfseries]
\textbf{Auspicio:} ACM SIGSOFT \& IEEE Computer Society (TCSE) \\
\textbf{Primera Edición:} 1975 \quad | \quad \textbf{Ranking CORE:} \textbf{A*} \quad | \quad \textbf{GGS Rating:} Clase 1 \\
\textbf{Índice $h5$ aproximado:} $\approx 88$ \quad | \quad \textbf{Tasa de Aceptación típica:} $18\% - 22\%$
\end{tcolorbox}
\textbf{Descripción e Importancia:} \\
ICSE es universalmente reconocida como la conferencia más prestigiosa y de mayor impacto en el campo de la ingeniería y desarrollo de software a nivel mundial. Cubre todos los aspectos del ciclo de vida del software, tales como ingeniería de requerimientos, diseño de arquitecturas, análisis estático y dinámico, pruebas automatizadas, ingeniería de software basada en inteligencia artificial (\textit{AI for SE}) y aspectos humanos en equipos de programadores. Publicar en ICSE constituye el mayor logro académico en la carrera de un investigador o laboratorio de software.

\vspace{0.4cm}

\begin{tcolorbox}[colback=white,colframe=tierAStar,title={\textbf{2. ACM ESEC/FSE -- Foundations of Software Engineering}},fonttitle=\bfseries]
\textbf{Auspicio:} ACM SIGSOFT \\
\textbf{Primera Edición:} 1993 \quad | \quad \textbf{Ranking CORE:} \textbf{A*} \quad | \quad \textbf{GGS Rating:} Clase 1 \\
\textbf{Índice $h5$ aproximado:} $\approx 68$ \quad | \quad \textbf{Tasa de Aceptación típica:} $20\% - 23\%$
\end{tcolorbox}
\textbf{Descripción e Importancia:} \\
ESEC/FSE es la conferencia emblemática europea/internacional de ACM dedicada a los fundamentos teóricos, matemáticos y empíricos del desarrollo de software. Es el foro donde se presentan modelos formales de concurrencia, síntesis de programas, verificación formal de código crítico, técnicas de depuración algorítmica y marcos de evolución de sistemas a gran escala.

\vspace{0.4cm}

\begin{tcolorbox}[colback=white,colframe=tierAStar,title={\textbf{3. ASE -- IEEE/ACM International Conference on Automated Software Engineering}},fonttitle=\bfseries]
\textbf{Auspicio:} IEEE Computer Society \& ACM SIGSOFT \\
\textbf{Primera Edición:} 1986 \quad | \quad \textbf{Ranking CORE:} \textbf{A*} \quad | \quad \textbf{GGS Rating:} Clase 1 \\
\textbf{Índice $h5$ aproximado:} $\approx 62$ \quad | \quad \textbf{Tasa de Aceptación típica:} $19\% - 21\%$
\end{tcolorbox}
\textbf{Descripción e Importancia:} \\
ASE se enfoca en la automatización integral de las tareas de desarrollo de software. En los últimos años ha cobrado un protagonismo excepcional al ser la sede principal de investigaciones pioneras en generación automática de código mediante modelos de lenguaje (LLMs), reparación automática de defectos (\textit{Automated Program Repair}), generación heurística de casos de prueba y asistentes de programación inteligentes.

\vspace{0.4cm}

\begin{tcolorbox}[colback=white,colframe=tierAStar,title={\textbf{4. ISSTA -- International Symposium on Software Testing and Analysis}},fonttitle=\bfseries]
\textbf{Auspicio:} ACM SIGSOFT \\
\textbf{Primera Edición:} 1991 \quad | \quad \textbf{Ranking CORE:} \textbf{A*} \quad | \quad \textbf{GGS Rating:} Clase 1 \\
\textbf{Índice $h5$ aproximado:} $\approx 51$ \quad | \quad \textbf{Tasa de Aceptación típica:} $21\% - 24\%$
\end{tcolorbox}
\textbf{Descripción e Importancia:} \\
ISSTA es el evento cumbre del mundo enfocado de manera exhaustiva en las fases de verificación, pruebas y análisis de software. Es la cuna de los métodos de \textit{fuzzing} moderno, ejecución simbólica, detección de vulnerabilidades en binarios y código fuente, y análisis dinámico de memoria para garantizar la fiabilidad del software crítico.

\vspace{0.4cm}

\begin{tcolorbox}[colback=white,colframe=tierAStar,title={\textbf{5. OOPSLA / SPLASH -- Object-Oriented Programming, Systems, Languages \& Applications}},fonttitle=\bfseries]
\textbf{Auspicio:} ACM SIGPLAN \\
\textbf{Primera Edición:} 1986 \quad | \quad \textbf{Ranking CORE:} \textbf{A*} \quad | \quad \textbf{GGS Rating:} Clase 1 \\
\textbf{Índice $h5$ aproximado:} $\approx 55$ \quad | \quad \textbf{Tasa de Aceptación típica:} $23\% - 26\%$
\end{tcolorbox}
\textbf{Descripción e Importancia:} \\
Perteneciente al macroevento SPLASH, OOPSLA conecta la teoría de los lenguajes de programación con la práctica del desarrollo de software orientado a objetos y multiparadigma. Aquí se han concebido conceptos revolucionarios de la industria, desde la programación orientada a aspectos, compiladores JIT, motores de recolección de basura modernos hasta frameworks de patrones de diseño.

% --- 3.2 NIVEL CORE A ---
\subsection{Nivel 2: Categoría CORE A (Alta Calidad y Reconocimiento Internacional)}

\begin{tcolorbox}[colback=white,colframe=tierA,title={\textbf{6. MSR -- International Conference on Mining Software Repositories}},fonttitle=\bfseries]
\textbf{Auspicio:} IEEE Computer Society \& ACM SIGSOFT \\
\textbf{Ranking CORE:} \textbf{A} \quad | \quad \textbf{Índice $h5$:} $\approx 52$ \quad | \quad \textbf{Tasa de Aceptación:} $25\% - 30\%$
\end{tcolorbox}
\textbf{Descripción:} MSR lidera la investigación empírica mediante la extracción masiva de datos en repositorios como GitHub, GitLab y Stack Overflow. Los artículos aquí presentados analizan patrones de commits, comportamiento de desarrolladores, análisis de dependencias de ecosistemas (npm, Maven, PyPI) y la evolución de la calidad del código a lo largo del tiempo.

\vspace{0.3cm}

\begin{tcolorbox}[colback=white,colframe=tierA,title={\textbf{7. ICSME -- IEEE International Conference on Software Maintenance and Evolution}},fonttitle=\bfseries]
\textbf{Auspicio:} IEEE Computer Society (TCSE) \\
\textbf{Ranking CORE:} \textbf{A} \quad | \quad \textbf{Índice $h5$:} $\approx 36$ \quad | \quad \textbf{Tasa de Aceptación:} $24\% - 28\%$
\end{tcolorbox}
\textbf{Descripción:} Conferencia primordial dedicada a una de las fases más costosas del ciclo de vida: el mantenimiento y evolución del software. Aborda la refactorización continua, gestión de deuda técnica, migración de sistemas monolíticos a microservicios y adaptación de sistemas heredados (\textit{legacy systems}).

\vspace{0.3cm}

\begin{tcolorbox}[colback=white,colframe=tierA,title={\textbf{8. ICST -- IEEE International Conference on Software Testing, Verification and Validation}},fonttitle=\bfseries]
\textbf{Auspicio:} IEEE Computer Society \\
\textbf{Ranking CORE:} \textbf{A} \quad | \quad \textbf{Índice $h5$:} $\approx 35$ \quad | \quad \textbf{Tasa de Aceptación:} $26\% - 30\%$
\end{tcolorbox}
\textbf{Descripción:} Foro líder impulsado por IEEE para investigadores y profesionales de pruebas de software. Presenta innovaciones en pruebas guiadas por modelos, pruebas de mutación, validación de software en sistemas embebidos y aseguramiento de calidad (\textit{QA}).

\vspace{0.3cm}

\begin{tcolorbox}[colback=white,colframe=tierA,title={\textbf{9. ESEM -- ACM/IEEE International Symposium on Empirical Software Engineering and Measurement}},fonttitle=\bfseries]
\textbf{Auspicio:} ACM \& IEEE Computer Society \\
\textbf{Ranking CORE:} \textbf{A} \quad | \quad \textbf{Índice $h5$:} $\approx 34$ \quad | \quad \textbf{Tasa de Aceptación:} $25\% - 29\%$
\end{tcolorbox}
\textbf{Descripción:} Conferencia rectora en la aplicación del método científico al desarrollo de software: estudios de caso industriales, encuestas a gran escala con desarrolladores, métricas de productividad humana y validación rigurosa de metodologías ágiles.

\vspace{0.3cm}

\begin{tcolorbox}[colback=white,colframe=tierA,title={\textbf{10. RE -- IEEE International Requirements Engineering Conference}},fonttitle=\bfseries]
\textbf{Auspicio:} IEEE Computer Society \\
\textbf{Ranking CORE:} \textbf{A} \quad | \quad \textbf{Índice $h5$:} $\approx 32$ \quad | \quad \textbf{Tasa de Aceptación:} $22\% - 26\%$
\end{tcolorbox}
\textbf{Descripción:} El evento más influyente dedicado a la elicitación, modelado, especificación y trazabilidad de requisitos de software, integrando hoy en día técnicas de procesamiento de lenguaje natural (NLP) para transformar requisitos de usuarios en especificaciones verificables.

\vspace{0.3cm}

\begin{tcolorbox}[colback=white,colframe=tierA,title={\textbf{11. MODELS -- ACM/IEEE International Conference on Model Driven Engineering}},fonttitle=\bfseries]
\textbf{Auspicio:} ACM \& IEEE Computer Society \\
\textbf{Ranking CORE:} \textbf{A} \quad | \quad \textbf{Índice $h5$:} $\approx 30$ \quad | \quad \textbf{Tasa de Aceptación:} $26\% - 31\%$
\end{tcolorbox}
\textbf{Descripción:} Centro neurálgico del desarrollo basado en modelos (MDD/MDE), lenguajes de dominio específico (DSL), metamodelado y arquitecturas generativas de software.

\vspace{0.3cm}

\begin{tcolorbox}[colback=white,colframe=tierA,title={\textbf{12. SANER -- IEEE International Conference on Software Analysis, Evolution and Reengineering}},fonttitle=\bfseries]
\textbf{Auspicio:} IEEE Computer Society \\
\textbf{Ranking CORE:} \textbf{A} \quad | \quad \textbf{Índice $h5$:} $\approx 33$ \quad | \quad \textbf{Tasa de Aceptación:} $26\% - 30\%$
\end{tcolorbox}
\textbf{Descripción:} Especializada en reingeniería de software, análisis estático de código a gran escala, comprensión de código fuente y deconstrucción de arquitecturas.

% --- 3.3 NIVEL CORE B ---
\subsection{Nivel 3: Categoría CORE B (Foros Especializados y Regionales de Prestigio)}
\begin{itemize}[leftmargin=*]
    \item \textbf{ICPC (IEEE/ACM International Conference on Program Comprehension):} Focalizada en comprender cómo los programadores leen, interpretan y modifican el código fuente, integrando estudios cognitivos y herramientas de visualización de software.
    \item \textbf{SEAMS (IEEE/ACM International Symposium on Software Engineering for Adaptive and Self-Managing Systems):} Orientada al diseño de software con capacidades de auto-adaptación, auto-reparación y auto-configuración en entornos dinámicos.
    \item \textbf{XP Conference (International Conference on Agile Software Development):} El foro académico-profesional pionero en métodos ágiles, Scrum, integración continua, desarrollo guiado por pruebas (TDD) y DevOps.
\end{itemize}

% ==============================================================================
% 4. MATRIZ COMPARATIVA DE RELEVANCIA
% ==============================================================================
\newpage
\section{Matriz Comparativa y Ordenamiento por Relevancia}

La siguiente tabla resume y ordena jerárquicamente las conferencias evaluadas, sintetizando su calificación en los principales índices internacionales:

\small
\renewcommand{\arraystretch}{1.25}
\begin{longtable}{@{} p{1.9cm} p{2.1cm} c c c c p{4.9cm} @{}}
\caption{Matriz general de conferencias de desarrollo de software ordenadas por relevancia académica.} \label{tab:conferencias} \\
\toprule
\rowcolor{tableHeader}
\textbf{Conferencia} & \textbf{Entidad} & \textbf{CORE} & \textbf{GGS} & \textbf{h5-idx} & \textbf{Tasa Acep.} & \textbf{Eje Temático Principal} \\
\midrule
\endfirsthead

\multicolumn{7}{c}%
{{\small\bfseries \tablename\ \thetable{} -- Continuación de la página anterior}} \\[0.2cm]
\toprule
\rowcolor{tableHeader}
\textbf{Conferencia} & \textbf{Entidad} & \textbf{CORE} & \textbf{GGS} & \textbf{h5-idx} & \textbf{Tasa Acep.} & \textbf{Eje Temático Principal} \\
\midrule
\endhead

\midrule
\multicolumn{7}{r@{}}{\small\textit{Continúa en la siguiente página \dots}} \\
\endfoot

\bottomrule
\endlastfoot

\rowcolor{tableStripe}
\textbf{ICSE} & ACM / IEEE & \textbf{A*} & Clase 1 & $\approx 88$ & $18\% - 22\%$ & Todo el espectro del ciclo de vida del software (Flagship mundial). \\
\textbf{ESEC/FSE} & ACM SIGSOFT & \textbf{A*} & Clase 1 & $\approx 68$ & $20\% - 23\%$ & Fundamentos teóricos, verificación, arquitectura y empirismo. \\
\rowcolor{tableStripe}
\textbf{ASE} & IEEE / ACM & \textbf{A*} & Clase 1 & $\approx 62$ & $19\% - 21\%$ & Automatización, síntesis de código, testing con IA y reparación. \\
\textbf{OOPSLA} & ACM SIGPLAN & \textbf{A*} & Clase 1 & $\approx 55$ & $23\% - 26\%$ & Lenguajes, paradigmas, compiladores y sistemas orientados a objetos. \\
\rowcolor{tableStripe}
\textbf{ISSTA} & ACM SIGSOFT & \textbf{A*} & Clase 1 & $\approx 51$ & $21\% - 24\%$ & Testing avanzado, fuzzing, análisis de código y detección de bugs. \\
\midrule
\textbf{MSR} & IEEE / ACM & \textbf{A} & Clase 2 & $\approx 52$ & $25\% - 30\%$ & Minería de datos en GitHub, analítica de código y telemetría. \\
\rowcolor{tableStripe}
\textbf{ICSME} & IEEE CS & \textbf{A} & Clase 2 & $\approx 36$ & $24\% - 28\%$ & Mantenimiento, evolución continua, refactorización y código legado. \\
\textbf{ICST} & IEEE CS & \textbf{A} & Clase 2 & $\approx 35$ & $26\% - 30\%$ & Pruebas en entornos industriales, verificación y validación formal. \\
\rowcolor{tableStripe}
\textbf{ESEM} & ACM / IEEE & \textbf{A} & Clase 2 & $\approx 34$ & $25\% - 29\%$ & Métricas de software, estudios empíricos y psicología de desarrollo. \\
\textbf{SANER} & IEEE CS & \textbf{A} & Clase 2 & $\approx 33$ & $26\% - 30\%$ & Reingeniería, ingeniería inversa y análisis de arquitecturas. \\
\rowcolor{tableStripe}
\textbf{RE} & IEEE CS & \textbf{A} & Clase 2 & $\approx 32$ & $22\% - 26\%$ & Elicitación, especificación y trazabilidad de requerimientos. \\
\textbf{MODELS} & ACM / IEEE & \textbf{A} & Clase 2 & $\approx 30$ & $26\% - 31\%$ & MDE, DSLs, diseño guiado por modelos y meta-arquitectura. \\
\midrule
\rowcolor{tableStripe}
\textbf{ICPC} & IEEE / ACM & \textbf{B} & Clase 2 & $\approx 28$ & $28\% - 32\%$ & Comprensión y cognición de código fuente por humanos y máquinas. \\
\textbf{SEAMS} & IEEE / ACM & \textbf{B} & Clase 2 & $\approx 24$ & $30\% - 35\%$ & Sistemas auto-adaptativos, robótica y software autonómico. \\
\rowcolor{tableStripe}
\textbf{XP Conf.} & Agile Alliance & \textbf{B} & Clase 2 & $\approx 22$ & $30\% - 35\%$ & Desarrollo ágil, Scrum, entrega continua y cultura DevOps. \\
\end{longtable}

% ==============================================================================
% 5. CONFERENCIAS INDUSTRIALES Y PROFESIONALES DE MAYOR IMPACTO
% ==============================================================================
\section{Conferencias Industriales y Profesionales de Mayor Impacto}

Paralelamente a los foros académicos arbitrados, la industria global del software cuenta con conferencias de difusión y networking profesional que definen la adopción tecnológica en empresas del sector productivo.

\subsection{Distinción entre Conferencias Académicas e Industriales}
\begin{itemize}[leftmargin=*]
    \item \textbf{Conferencias Académicas (ICSE, FSE, ASE):} Buscan la \textit{novedad científica}, el rigor metodológico, la generalización de conclusiones y la reproducibilidad mediante pruebas empíricas. Se enfocan en el horizonte de los próximos 3 a 7 años.
    \item \textbf{Conferencias Industriales (QCon, KubeCon, Devoxx):} Se enfocan en la \textit{aplicabilidad inmediata}, lecciones aprendidas en producción a escala masiva (\textit{war stories}), resiliencia operativa y adopción de stacks modernos.
\end{itemize}

\subsection{Las Conferencias Industriales Más Relevantes}
\begin{enumerate}[leftmargin=*]
    \item \textbf{QCon (organizada por InfoQ):}
    Considerada la conferencia líder mundial para arquitectos de software, ingenieros sénior y líderes de equipo. Se enfoca en patrones de diseño en microservicios, sistemas distribuidos, plataformas en la nube y cultura técnica.
    
    \item \textbf{KubeCon + CloudNativeCon (organizada por la CNCF / Linux Foundation):}
    El evento más concurrido a nivel global sobre computación nativa en la nube, contenedores (Docker, Kubernetes), observabilidad, mallas de servicios (\textit{service mesh}) y \textit{Platform Engineering}.
    
    \item \textbf{Strange Loop / Systems Distributed:}
    Reconocida por su excelencia técnica en programación funcional, bases de datos distribuidas y lenguajes experimentales. Reúne a la vanguardia de desarrolladores de sistemas de alto rendimiento.
    
    \item \textbf{GOTO Conferences:}
    Conferencias orientadas a las mejores prácticas de artesanía de software (\textit{software craftsmanship}), refactorización, \textit{Domain-Driven Design (DDD)} y metodologías ágiles pragmáticas.
    
    \item \textbf{Devoxx:}
    La conferencia comunitaria más grande para desarrolladores de la plataforma Java, ecosistema JVM, Spring Boot, microservicios y desarrollo backend.
    
    \item \textbf{Conferencias Flagship de Grandes Proveedores Tecnológicos:}
    \begin{itemize}
        \item \textbf{AWS re:Invent:} Anuncios clave sobre infraestructura en la nube, bases de datos y servicios de cómputo serverless.
        \item \textbf{Google I/O:} Plataformas de desarrollo móvil (Android, Flutter), Web y modelos fundacionales de Inteligencia Artificial.
        \item \textbf{Microsoft Build:} Herramientas para desarrolladores (Visual Studio, GitHub, Azure, .NET) y copilotos de IA.
        \item \textbf{Apple WWDC:} Novedades en lenguajes como Swift, arquitectura de sistemas y ecosistemas móviles / de escritorio.
    \end{itemize}
\end{enumerate}

% ==============================================================================
% 6. TENDENCIAS EMERGENTES EN LAS CONFERENCIAS ACTUALES
% ==============================================================================
\section{Tendencias Actuales de Investigación (2024 -- 2026)}

Al examinar los programas técnicos recientes de las conferencias A* (ICSE, FSE, ASE), se evidencian cuatro corrientes temáticas dominantes:

\begin{enumerate}[leftmargin=*]
    \item \textbf{AI for Software Engineering (AI4SE):} Aplicación de modelos de lenguaje masivos (LLMs) y agentes autónomos para la síntesis de código, generación automática de casos de prueba unitarios, detección de brechas de seguridad y refactorización inteligente.
    \item \textbf{Software Engineering for AI (SE4AI):} Creación de metodologías formales para probar, auditar, versionar y monitorear sistemas que incorporan modelos no deterministas y aprendizaje automático (\textit{MLOps}).
    \item \textbf{Seguridad en la Cadena de Suministro de Software (Software Supply Chain Security):} Generación y verificación de listas de materiales de software (\textit{SBOM}), protección contra vulnerabilidades en dependencias transitorias y compilación verificable (\textit{SLSA framework}).
    \item \textbf{Green Software e Ingeniería Sostenible:} Métricas y optimizaciones para reducir la huella de carbono y el consumo de energía en centros de datos derivado de algoritmos ineficientes.
\end{enumerate}

% ==============================================================================
% 7. CONCLUSIONES
% ==============================================================================
\section{Conclusiones}

\begin{itemize}[leftmargin=*]
    \item \textbf{Jerarquía y Estándares de Excelencia:} En el campo del desarrollo de software, las conferencias con calificación \textbf{CORE A*} (principalmente \textbf{ICSE}, \textbf{FSE}, \textbf{ASE}, \textbf{OOPSLA} e \textbf{ISSTA}) definen la cúspide del conocimiento formal, alcanzando índices de impacto comparables o superiores a los mejores journals de la disciplina.
    \item \textbf{Especialización Temática:} El ordenamiento por relevancia evidencia una sólida estructuración por áreas de especialidad: mientras que ICSE aborda el ciclo de vida completo, foros como MSR, ICSME e ISSTA lideran con maestría subdominios críticos como la analítica de repositorios, el mantenimiento de sistemas y el aseguramiento de la calidad.
    \item \textbf{Complementariedad con la Práctica Profesional:} Existe una retroalimentación continua entre las conferencias científicas y las de la industria (como QCon y KubeCon). Los descubrimientos algorítmicos concebidos en las conferencias académicas suelen convertirse en herramientas y estándares comerciales en un plazo de 3 a 5 años.
    \item \textbf{Relevancia para la Formación en la UNSAAC:} Para los estudiantes y futuros tesistas de la Escuela Profesional de Ingeniería Informática y de Sistemas de la UNSAAC, la consulta directa de las actas de estas conferencias (disponibles en la \textit{ACM Digital Library} y \textit{IEEE Xplore}) representa una oportunidad invaluable para identificar el estado del arte y fundamentar proyectos de investigación formativa de alto nivel.
\end{itemize}

% ==============================================================================
% REFERENCIAS BIBLIOGRÁFICAS
% ==============================================================================
\newpage
\begin{thebibliography}{99}

\bibitem{core2023}
Computing Research and Education Association of Australasia (CORE). (2023). 
\textit{CORE Conference Rankings Portal}. 
Disponible en: \url{http://portal.core.edu.au/conf-ranks/}

\bibitem{googlescholar}
Google Scholar Metrics. (2025). 
\textit{Top Publications -- Software Systems}. 
Disponible en: \url{https://scholar.google.com/citations?view_op=top_venues&hl=es&vq=eng_softwaresystems}

\bibitem{csrankings}
Berger, E. (2024). 
\textit{CSRankings: Computer Science Rankings}. 
Disponible en: \url{https://csrankings.org/}

\bibitem{ieee_tcse}
IEEE Computer Society Technical Council on Software Engineering (TCSE). (2025). 
\textit{Sponsored and Co-Sponsored Conferences: ICSE, ICSME, ASE, SANER}. 
Disponible en: \url{https://www.computer.org/}

\bibitem{acm_sigsoft}
ACM Special Interest Group on Software Engineering (SIGSOFT). (2025). 
\textit{Conferences and Events: ICSE, FSE, ISSTA, ESEM}. 
Disponible en: \url{https://www.sigsoft.org/}

\bibitem{sommerville}
Sommerville, I. (2016). 
\textit{Software Engineering} (10th ed.). Pearson Education.

\bibitem{pressman}
Pressman, R. S., \& Maxim, B. R. (2020). 
\textit{Software Engineering: A Practitioner's Approach} (9th ed.). McGraw-Hill Education.

\end{thebibliography}

\end{document}
